\documentclass{beamer}

\usepackage{graphicx}
\usepackage{tikz}
\usetikzlibrary{external}
\usepackage[mode=buildnew]{standalone} % Allows inclusion of standalone files

\usepackage{caption}
\usepackage{subcaption}



\usetheme{default} % Change "default" to a specific theme, e.g., "Madrid" or "Berlin"
\title{ORCA: Observability with Realtime Control and Aggregation}
\author[Left Text]{Ankush Jain}
\institute{CMU}
\date{\today}

\begin{document}

\begin{frame}
  \titlepage % Generates the title slide
\end{frame}

\begin{frame}{ORCA: Key Ideas}
  ORCA is designed to answer any observability query using a combination of two capabilities in a dynamic, online, realtime manner.

  \begin{enumerate}
    \item Control over:
      \begin{itemize}
        \item What telemetry streams are active.
        \item What probes in each stream are active.
      \end{itemize}
    \item Control over what aggregations are applied to that telemetry. ORCA automatically orchestrates those aggregations in the best possible manner.
  \end{enumerate}
\end{frame}

\begin{frame}{ORCA: Workflow}
  \begin{itemize}
  \item An application can always afford to run under ORCA. In the common case, no telemetry is published, and no overhead is incurred.

  \item If a user chooses to enable a probe, and an aggregation for that probe, a dataflow for that aggregation is orchestrated using SQL operators + Arrow.

  \item The aggregation transfers the least amount of data necessary across nodes. The bulk of the aggregation is done in dedicated aggregators that are like MapReduce nodes running SQL partitioning/shuffling/aggregation. Preliminary filtering and aggregation can be done on leaf nodes using operator pushdown. This is very standard SQL query planning.
\end{itemize}
\end{frame}

\begin{frame}{ORCA: Concepts}
  A \textbf{stream} represents a schema. A \textbf{probe} represents a source generating something that fits the schema.

\begin{table}[h]
    \centering
    \begin{tabular}{|l|l|l|l|l|}
        \hline
        \textbf{Timestamp} & \textbf{Timestep} & \textbf{Rank} & \textbf{Probe ID} & \textbf{Probe Value} \\ \hline
        1234567890         & 10                & 0             & 101               & 42.0                 \\ \hline
        1234567891         & 10                & 1             & 102               & 36.5                 \\ \hline
        1234567892         & 11                & 0             & 103               & 28.4                 \\ \hline
        1234567893         & 11                & 1             & 104               & 22.1                 \\ \hline
    \end{tabular}
    \caption{Per-timestep Telemetry Dataframe}
    \label{tab:telemetry-data}
\end{table}

\textbf{Usecases}: say you want to monitor an AMR application. You may define a schema for block to rank mapping, and another schema for blockwise telemetry. Your aggregation query may want to join the two.
\end{frame}

\begin{frame}{ORCA: Components}
  \begin{itemize}
    \item \textbf{Controller}: Triggers orchestration, receives aggregates
    \item \textbf{Aggregators}: On-fabric nodes that form an overlay aggregation network. They run executors that compute per-partition shards of a query.
    \item \textbf{Parallel app 3-hop}: The application itself can assist the workflow by 1) broadcasting control messages via an internal overlay, and 2) doing leaf aggregations---preaggregation, prefiltering etc.
  \end{itemize}
\end{frame}

% \begin{frame}{Control Path}
% \begin{figure}[ht]
%     \centering
%     \includestandalone[width=0.8\linewidth]{../tikz-exps/overlay-ctrl} % Embeds the standalone file
%     \caption{Overlay - Control Path}
%     \label{fig:overlay_ctrl}
% \end{figure}
% \end{frame}
%
% \begin{frame}{Aggregation Path}
% \begin{figure}[ht]
%     \centering
%     \includestandalone[width=0.8\linewidth]{../tikz-exps/overlay-agg-new} % Embeds the standalone file
%     \caption{Overlay - Control Path}
%     \label{fig:overlay_data}
% \end{figure}
% \end{frame}
%
\begin{frame}{ORCA: Stack}
  \begin{itemize}
    \item Query planning using \textbf{datafusion}, a rust-based query planning engine from the Apache foundation.
    \item \textbf{Arrow} for all in-memory telemetry. \textbf{Parquet} maybe for on-disk telemetry (can be set to persist conditionally).
    \item \textbf{Mercury} + QoS for all communication.
  \end{itemize}

Bulk of the codebase is still going to be C/C++. Only the query engine bits will be Rust (Rust-C FFI + cbindgen).
\end{frame}


\begin{frame}[fragile]{Datafusion/Orca: Query Planning}

\textbf{SQL}
  \begin{verbatim}
SELECT max(time)
FROM table
GROUP BY event_id;
  \end{verbatim}

\textbf{Logical Plan}
  \begin{verbatim}
Projection: max(time)
  Aggregate: groupBy=[event_id], aggr=[max(time)]
    TableScan: table, projection=[event_id, time]
  \end{verbatim}
\end{frame}

\begin{frame}[fragile]{Datafusion/Orca: Query Planning II}
\textbf{Physical Plan}
    \begin{verbatim}
Projection: max(time)
  HashAggregate: groupBy=[event_id], aggr=[max(time)]
    Repartition: hash(event_id)
      TableScan: table, projection=[event_id, time]
    \end{verbatim}

  \textbf{Distributed Query Plan}
  \begin{verbatim}
Stage 2:
ShuffleWrite: partitionBy=[event_id]
  PartialAggregate: groupBy=[event_id], aggr=[max(time)]
    TableScan: table, projection=[event_id, time]

Stage 1:
Projection: max(time)
  FinalAggregate: groupBy=[event_id], aggr=[max(time)]
    ShuffleRead: partitionBy=[event_id]
  \end{verbatim}
  % \end{block}
\end{frame}

\begin{frame}{ORCA: Misc Details}
  \begin{itemize}
    \item A telemetry dataframe is pushed by each rank at the end of a \emph{timestep}. It contains all data collected within that timeframe.
    \item The most common telemetry aggregation is envisioned to be per-timestep. All data for a timestep is collected, and an online OLAP query is run. Thus ORCA is not really streaming. In SQL parlance, there is a window and data is batched as per that window. This also enables more efficient use of the fabric.
    \item Control message broadcast will use timestep-linked two-phase commit for consistency. It can either be synchronous or async --- that is a performance decision. It will always be applied at the same timestep on all ranks.
  \end{itemize}
\end{frame}

\begin{frame}{ORCA: Extension Potential}
  I think ORCA has a lot of extension potential. Idea: orchestrate arbitrary interactions with a bulk-synchronous application.

  \begin{itemize}
    \item Automatically monitor PMPI/Kokkos calls using their respective profiling interfaces.
    \item Applications can declare specific inspection point for ORCA to access (per-block compute kernel costs etc.)
    \item Tuning knobs can be exposed to Orca for online tuning (Control Information Base).
    \item Controller can even participate in placement, steering etc. through the control interface.
    \item eBPF extensions for uprobe/kprobe data.
    \item Reuse infra to shuffle, index, and aggregate scientific data. Particle data is already Arrow-friendly. Extending SQL OLAP to multidimensional mesh data remains a possibility (\texttt{arrow-zarr}).
  \end{itemize}
\end{frame}

\begin{frame}{ORCA: Misc Things}
  \begin{itemize}
    \item Datafusion Catalog API to offer and discover streams
    \item Accelerated operators (or DPU operators) can be implemented transparently for any stage of the query plan.
    \item Sketch-based approximate implementations for certain query types possible (work ongoing at CMU).
    \item Other OLAP world developments: Substrait, ADBC --- may be relevant.
    \item Query language syntax is orthogonal. Can be SQL, SQL++, or something else entirely. Datafusion allows custom parsers.
    \item Datafusion offers single-node query planning. Distributed query planning is additional code on top. Largely exists in \texttt{datafusion-ray}.
  \end{itemize}
\end{frame}


\end{document}
